% Long-time distribution targets; detailed proof follows at the end of this structure.

有限观测时间不会产生严格的能量守恒 $\delta$ 函数，而产生宽度约为 $\hbar/T$ 的有限时间谱核。令 $\delta E$ 为实能量失配，单位 J；令 $T>0$ 为相互作用或观测时间，单位 s。定义
\begin{equation*}
 F_T(\delta E)
 \equiv T\,\operatorname{sinc}^2\!\left(\frac{\delta E\,T}{2\hbar}\right),
 \qquad
 \operatorname{sinc}x\equiv\frac{\sin x}{x}.
\end{equation*}
因为 $F_T$ 的峰高随 $T$ 增长而峰宽随 $1/T$ 缩小，它不能按普通逐点函数极限理解；正确极限必须用分布作用在测试函数上的方式定义。

对任意光滑、足够快衰减的测试函数 $g(\delta E)$，定义
\begin{equation*}
 I_T[g]\equiv
 \int_{-\infty}^{\infty}\dd(\delta E)\,
 F_T(\delta E)g(\delta E).
\end{equation*}
作无量纲变量替换
$x=\delta E\,T/(2\hbar)$，于是
\begin{equation}
 \boxed{
 I_T[g]
 =2\hbar\int_{-\infty}^{\infty}\dd x\,
 \left(\frac{\sin x}{x}\right)^2
 g\!\left(\frac{2\hbar x}{T}\right).}
 \label{eq:sinc-test-substitution-r15}
\end{equation}
因此长时间极限只剩两个问题：$g(2\hbar x/T)\to g(0)$，以及核 $\operatorname{sinc}^2x$ 的总面积是多少。

为不把 Dirichlet 积分作为未证公式，引入指数调节器 $\epsilon>0$，定义
\begin{equation*}
 D_\epsilon(b)
 \equiv\int_0^\infty\dd x\,\ee^{-\epsilon x}\frac{\sin bx}{x}.
\end{equation*}
对 $b$ 求导后积分变成初等 Laplace 变换，最终得到
\begin{equation}
 \boxed{
 D_\epsilon(b)=\arctan\frac{b}{\epsilon}
 \xrightarrow[\epsilon\to0^+]{}
 \frac{\pi}{2}\operatorname{sgn}b.}
 \label{eq:dirichlet-integral-derived-r15}
\end{equation}
由此进一步得到
$\int_{-\infty}^{\infty}(\sin x/x)^2\dd x=\pi$。代回\eqnrefs{eq:sinc-test-substitution-r15} 后，有限时间谱核的分布极限为
\begin{equation}
 \boxed{
 \lim_{T\to\infty}
 T\,\operatorname{sinc}^2\!\left(\frac{\delta E\,T}{2\hbar}\right)
 =2\pi\hbar\,\delta(\delta E).}
 \label{eq:sinc-delta-proved-r15}
\end{equation}
\eqnrefs{eq:sinc-delta-proved-r15} 的等号仅表示对测试函数积分后的分布收敛；任何有限 $T$ 的实际过程仍必须保留宽度有限的 $\operatorname{sinc}^2$ 线形，不能把它提前替换成严格 $\delta$ 函数。
